\documentclass[11pt,a4paper,twoside]{article}
\usepackage[utf8]{inputenc}
\usepackage[english]{babel}
\usepackage{actuarialangle, actuarialsymbol}

\begin{document}
Define comonotonicity and Denneberg's equivalences

Two random variable \( X, Y : \Omega \to \mathbb{R} \) are comonotone if
\( \forall \omega, \omega' \in \Omega \),
\( (X(\omega) - X(\omega'))(Y(\omega) - Y(\omega')) \geq 0\)

The following assertions are equivalent for random variables \(X, Y\):

\(X\) and \(Y\) are comonotone.
\(\exists\) random variable \(Z\) and nondecreasing \(u, v: \mathbb{R} \rightarrow \mathbb{R}\) s.t. \(X = u(Z)\) and \(Y = v(Z)\).
\(\exists\) nondecreasing \(u, v \in C(\mathbb{R}), u + v = id_{\mathbb{R}}\) s.t. for \(Z := X + Y, X = u(Z), Y = v(Z)\).

Note that \(u, v\) in condition 3. are necessarily 1-Lipschitz-continuous, i.e. \(|u(x) - u(y)| \le |x - y|\) and similarly for \(v\).


State Schmeidler's theorem

Assume \(\varphi: L^\infty \rightarrow \mathbb{R}\) has the following properties:

\(|\cdot|_\infty\)-continuity
\(\varphi(1_\Omega) = 1\)
monotone: \(X(w) \le Y(w)\) for all \(w \in \Omega \implies \varphi(X) \le \varphi(Y)\)

Set \(\mu(A) := \varphi(1_A), A \in \mathcal{F}\) (see later: \(\mu\) Choquet capacity). Then the following are equivalent:

1. \(\varphi\) is comonotonic additive: \(X, Y \in L^\infty\) comonotone \(\implies \varphi(X + Y) = \varphi(X) + \varphi(Y)\).
2. \(\varphi = \int \cdot d \mu\) (meaning its a Choquet integral)


State 4 axioms for premium principles (in the context of Choquet integrals)

\(\mathcal{X} \hat{=}\) nonnegative r.v.s on probability space \((\Omega, \mathcal{F}, \mathbb{P})\)

Authors posit four axioms for premium principle \(\pi: \mathcal{X} \rightarrow [0, \infty]\) with \(\pi(1_\Omega) = 1\):

A1 Law-invariance: If \(X, Y \in \mathcal{X}\) satisfy \(\mathbb{P} \circ X^{-1} = \mathbb{P} \circ Y^{-1}\), then \(\pi(X) = \pi(Y)\)
A2 Monotonicity: \(X(w) \le Y(w)\) for all \(w \in \Omega \implies \pi(X) \le \pi(Y)\)
A3 Comonotonic additivity: \(X, Y \in \mathcal{X}\) comonotone \(\implies \pi(X + Y) = \pi(X) + \pi(Y)\)
A4 Continuity: \(\forall X \in \mathcal{X}: \lim_{d \downarrow 0} \pi((X - d)^+) = \pi(X) = \lim_{d \uparrow \infty} \pi(X \wedge d)\)

Technical condition on \((\Omega, \mathcal{F}, \mathbb{P}) \implies \exists_1\) concave distortion function \(h: [0, 1] \rightarrow [0, 1]\)

s.t.
\(\pi(X) = \int X d(h \circ \mathbb{P}), X \in \mathcal{X}\)


Define distorted probability prices and the related proposition
for their premiums

Let \((\Omega, \mathcal{F}, \mathbb{P})\) be a probability space.

A continuous (non-decreasing) concave function \(h:[0,1]\rightarrow[0,1]\) with \(h(0)=0\) and
\(h(1)=1\) is a concave distortion function.
The h-distorted probability price of claim \(S\ge0\) is
\(\pi_{h}(S):=\int S~d(h\circ\mathbb{P})=\int_{0}^{\infty}h(\mathbb{P}(S>x))dx.\)

Concavity \(\rightarrow\) more weight on events with small \(\mathbb{P}\)-probability

\(h\) concave distortion function as above.
\(S, S'\) nonnegative random variables.
\(t > 0, 0 < \lambda < 1\).

Then:

\(h(p) \ge p, p \in [0, 1]\)
\(\pi_h\) is positively homogeneous.
\(\pi_h(S) \ge \mathbb{E}[S]\)
If \(h(p) > p, 0 < p < 1\), and \(S\) not deterministic, then \(\pi_h(S) > \mathbb{E}[S]\)
\((\Omega, \mathcal{F}, \mathbb{P})\) atomless, then \(\pi_h\) is convex:

\(\pi_h(\lambda S + (1 - \lambda)S') \le \lambda \pi_h(S) + (1 - \lambda)\pi_h(S').\)

\end{document}
